\documentclass[11pt]{article}
\usepackage[margin=0.88in]{geometry}
\usepackage{amsmath,amssymb,amsthm,mathtools,microtype,graphicx}
\usepackage[hidelinks]{hyperref}

\newtheorem{theorem}{Theorem}
\newtheorem{lemma}{Lemma}
\newcommand{\D}{\mathbb D}
\newcommand{\T}{\mathbb T}
\newcommand{\C}{\mathbb C}
\DeclareMathOperator{\spec}{spec}

\title{Disconnected spectra for constrained Toeplitz operators}
\author{Candidate counterexamples and an essential-spectrum theorem\\for Question 6.1 of arXiv:2110.06379}
\date{17 August 2026}

\begin{document}
\maketitle

\begin{center}
\fbox{\parbox{0.93\textwidth}{\textbf{Answer to two parts of Question 6.1.}
Ordinary spectral connectedness is false in both the Neil and two-point
settings, even for the bounded analytic rational symbol
\[
                  \phi(z)=\frac{z}{(1-z/2)^2}.
\]
In contrast, for every $L^\infty$ symbol and every finite-codimensional
Hardy subspace, the constrained compression has exactly the same Fredholm
essential spectrum as the classical Toeplitz operator.  Its essential
spectrum is therefore always connected.}}
\end{center}

\section{The source question}

Felder, Pfeffer, and Russo \cite{FPR} consider compressions
\[
        T_\phi^M=P_M(\phi\,\cdot)\big|_M
\]
to representation-carrying finite-codimensional subspaces $M\subset H^2$.
Their Question 6.1, on source PDF page 17, asks whether the spectrum,
essential spectrum, and relative point spectrum are connected for arbitrary
$\phi\in L^\infty$ in the two-point and Neil settings.  Their Question 6.2
asks the relative-spectrum question for general finite-codimensional
subalgebras.

\begin{figure}[h]
\centering
\includegraphics[width=0.90\textwidth]{figures/open_problem_crop.png}
\caption{Questions 6.1 and 6.2 in the source PDF, page 17.}
\end{figure}

\section{Idea of the proof}

Each Neil or two-point Hardy space in the examples splits as
\[
                   M=\C f\oplus \Theta H^2,
\]
where $\Theta$ is a degree-two finite Blaschke product.  An analytic symbol
preserves $\Theta H^2$.  Relative to this splitting, its constrained
Toeplitz operator is lower triangular:
\[
 T_\phi^M\cong
 \begin{pmatrix}
   \gamma&0\\ *&M_\phi
 \end{pmatrix},
 \qquad
 \gamma=\frac{\langle\phi f,f\rangle}{\|f\|^2}.
\]
Thus any $\gamma$ outside $\phi(\overline\D)$ becomes an isolated spectral
point.  The chosen rational $\phi$ and explicit $f$ make this happen in both
settings.

The essential-spectrum statement is even simpler.  Extend the compression
by zero on the finite-dimensional orthogonal complement of $M$.  The
extended operator differs from the classical $T_\phi$ by finite rank.
Fredholm essential spectra are invariant under that operation and under
adding a finite-dimensional direct summand.

\section{A triangular compression lemma}

\begin{lemma}\label{lem:triangle}
Let $\Theta$ be a finite Blaschke product, let $0\ne f\in
K_\Theta:=H^2\ominus\Theta H^2$, and suppose
$M=\C f\oplus\Theta H^2$ orthogonally.  If $\phi$ is analytic on a
neighborhood of $\overline\D$, put
\[
                  \gamma_f=\frac{\langle\phi f,f\rangle}{\|f\|^2}.
\]
If $\gamma_f\notin\phi(\overline\D)$, then
\[
                 \sigma(T_\phi^M)
                 =\{\gamma_f\}\cup\phi(\overline\D).
\]
In particular, this spectrum is disconnected and $\gamma_f$ is an isolated
eigenvalue.
\end{lemma}

\begin{proof}
Multiplication by $\phi$ preserves $\Theta H^2$.  After normalizing $f$ and
identifying $\Theta H^2$ unitarily with $H^2$, the compression has the block
form
\[
 T_\phi^M=
 \begin{pmatrix}\gamma_f&0\\ g&M_\phi\end{pmatrix}
\]
for some $g\in H^2$.  Since $\phi$ is analytic past the closed disk,
$M_\phi=\phi(M_z)$ and spectral mapping gives
$\sigma(M_\phi)=\phi(\overline\D)$.

For $\lambda\ne\gamma_f$, elementary block factorization shows that
$T_\phi^M-\lambda$ is invertible exactly when $M_\phi-\lambda$ is
invertible.  At $\lambda=\gamma_f$, the lower-right block is invertible by
hypothesis while the scalar diagonal block vanishes; hence the full operator
is not invertible.  This proves the displayed spectrum.  Because
$\phi(\overline\D)$ is compact and does not contain $\gamma_f$, the singleton
is an isolated connected component.  Solving the lower block with
$(M_\phi-\gamma_f)^{-1}$ also gives a nonzero eigenvector.
\end{proof}

\section{Neil-algebra counterexample}

Consider the Neil representation space
\[
 M_N=\left\{h\in H^2:h(0)=-h'(0)\right\}
     =\C f_N\oplus z^2H^2,
 \qquad f_N(z)=\frac{1-z}{\sqrt2}.
\]
This corresponds to the source parameters
$(\alpha,\beta)=(2^{-1/2},-2^{-1/2})$.  Take
\[
                         \phi(z)=\frac{z}{(1-z/2)^2}.
\]
The symbol is bounded and analytic on a neighborhood of $\overline\D$.
Since its constant and linear Taylor coefficients are $0$ and $1$,
respectively,
\[
 \gamma_N=\langle\phi f_N,f_N\rangle=-\frac12.
\]
Moreover,
\[
 \phi(z)=-\frac12
 \quad\Longleftrightarrow\quad
 2z=-(1-z/2)^2
 \quad\Longleftrightarrow\quad
 (z+2)^2=0.
\]
The only preimage is $z=-2$, so $-1/2\notin\phi(\overline\D)$.
Lemma \ref{lem:triangle} gives the exact disconnected spectrum
\[
             \boxed{\displaystyle
             \sigma(T_\phi^{M_N})
             =\left\{-\frac12\right\}\cup\phi(\overline\D).}
\]

\section{Two-point counterexample}

Now set
\[
 a=-\frac1{10},\qquad b=\frac1{10},\qquad t=\frac{11}{9},
\]
and let $\Theta=B_{a,b}$.  The function
\[
              f(z)=\frac{1-z}{1-z^2/100}
\]
belongs to $K_\Theta$.  Direct evaluation gives
\[
 \frac{f(a)}{f(b)}=\frac{1+1/10}{1-1/10}=\frac{11}{9}=t.
\]
Consequently the two-point Hardy space in the source has the orthogonal
decomposition
\[
                    H_t^2=\C f\oplus B_{a,b}H^2.
\]

We compute its scalar block exactly.  Put $q=1/100$.  The Taylor
coefficients of $f$ occur in pairs,
\[
              f_{2j}=q^j,\qquad f_{2j+1}=-q^j,
\]
and $\|f\|^2=2/(1-q^2)$.  Therefore the normalized correlations are
\[
 \frac{\langle z^{2\ell}f,f\rangle}{\|f\|^2}=q^\ell,\qquad
 \frac{\langle z^{2\ell+1}f,f\rangle}{\|f\|^2}
 =-\frac{1+q}{2}q^\ell.
\]
Since
\[
                 \phi(z)=\sum_{k\ge1}\frac{k}{2^{k-1}}z^k,
\]
putting $x=q/4$ and summing the two geometric derivatives yields
\begin{align*}
 \gamma_t
 &=4\sum_{\ell\ge1}\ell x^\ell
   -\frac{1+q}{2}\sum_{\ell\ge0}(2\ell+1)x^\ell\\
 &=\frac{4x}{(1-x)^2}
   -\frac{(1+q)(1+x)}{2(1-x)^2}
 =-\frac{79402}{159201}.
\end{align*}
Write $s=79402/159201$.  Exact arithmetic gives
\[
                         \frac49<s<\frac12.
\]
The equation $\phi(z)=-s$ is
\[
                  \frac{s}{4}z^2+(1-s)z+s=0.
\]
Its two roots are real and negative.  The smaller absolute value is
\[
       \frac{2((1-s)-\sqrt{1-2s})}{s},
\]
which exceeds $1$ exactly when $s>4/9$.  Hence neither root lies in
$\overline\D$, and $\gamma_t\notin\phi(\overline\D)$.  Lemma
\ref{lem:triangle} now gives a second exact disconnected spectrum:
\[
       \boxed{\displaystyle
       \sigma(T_\phi^t)
       =\left\{-\frac{79402}{159201}\right\}
        \cup\phi(\overline\D).}
\]

\section{Essential spectra remain connected}

\begin{theorem}\label{thm:essential}
Let $M\subset H^2$ be any closed finite-codimensional subspace and
$\phi\in L^\infty(\T)$.  Then the Fredholm essential spectrum of the
compression satisfies
\[
                         \sigma_{\mathrm e}(T_\phi^M)
                         =\sigma_{\mathrm e}(T_\phi).
\]
Consequently every associated compression in Question 6.1, and more
generally every finite-codimensional Hardy compression, has connected
essential spectrum.
\end{theorem}

\begin{proof}
Let $F=H^2\ominus M$, which is finite-dimensional, and extend the compression
to $H^2=M\oplus F$ by
\[
                    \widetilde T=T_\phi^M\oplus0_F.
\]
For $m\in M$ and $f\in F$,
\[
 (T_\phi-\widetilde T)(m+f)
 =T_\phi f+(I-P_M)T_\phi m.
\]
The first term ranges in the finite-dimensional space $T_\phi F$ and the
second in $F$.  Thus $T_\phi-\widetilde T$ has finite rank, and Atkinson's
theorem gives
\[
                    \sigma_{\mathrm e}(\widetilde T)
                    =\sigma_{\mathrm e}(T_\phi).
\]
Adding or deleting the finite-dimensional zero summand does not affect
Fredholmness (including at $\lambda=0$), so
$\sigma_{\mathrm e}(\widetilde T)=\sigma_{\mathrm e}(T_\phi^M)$.
The classical essential spectrum on the right is connected by the
Widom--Douglas theorem quoted in the source's introduction.
\end{proof}

\section{Verification, novelty, and remaining scope}

The included exact-arithmetic script checks the Neil factorization, the
two-point correlation sum, the rational value
$-79402/159201$, both strict bounds on $s$, both preimages outside the unit
disk, $t=11/9$, and $\|f\|^2=20000/9999$.  No computation is used as a proof.

A bounded novelty search on 17 August 2026 covered the run's four lightweight
indexes and web searches using the exact title, arXiv id, Question 6.1, the
phrase \emph{disconnected spectrum}, \emph{Neil algebra}, and
finite-codimensional Toeplitz compressions.  The exact hits were the source
and its published version; no later explicit resolution or matching
counterexample was found.  Novelty is plausible but not certified.

These counterexamples completely answer the ordinary-spectrum part of
Question 6.1 in the negative for both named settings.  Theorem
\ref{thm:essential} completely answers its essential-spectrum part in the
positive, in greater generality.  The relative point-spectrum part of
Question 6.1 and Question 6.2 remain open here.  Human review should focus on
the block decomposition and the exact two-point correlation calculation.
Verdict: \textbf{candidate counterexamples and full essential-spectrum
theorem, likely valid}.

\begin{thebibliography}{9}
\bibitem{FPR}
C. Felder, D. T. Pfeffer, and B. P. Russo,
\emph{Spectra for Toeplitz Operators Associated with a Constrained
Subalgebra}, arXiv:2110.06379; Integral Equations Operator Theory 94 (2022),
Article 15.
\end{thebibliography}

\end{document}
